% Revize 2026-10-28: obsah, matematika a návaznost; podrobnosti a časy
% viz planovani/REVIZE_PRIPRAV_1_ROCNIK.md.
% Aktuální porovnání se skutečným obsahem místních PDF Realisticky.cz:
% externi_zdroje/realisticky/01_Zakladni_poznatky/02_Ciselne_obory/010205_Realna_cisla_I/010205_Lekce.pdf
% externi_zdroje/realisticky/01_Zakladni_poznatky/02_Ciselne_obory/010206_Realna_cisla_II/010206_Lekce.pdf
% externi_zdroje/realisticky/01_Zakladni_poznatky/02_Ciselne_obory/010207_Druha_odmocnina/010207_Lekce.pdf
% Rozpis: 1_rocnik/tematicke_plany/Rozpis_vyuky_MAT_1_rocnik.tex, hodina 7;
% SPS.pdf, oddíl 1.2. Návaznost: 06 (zlomky), 08 (intervaly), 12 (konvence oborů).
% Učebnice: ucebnice/01_Zakladni_poznatky/Ciselne_obory/Číselné obory.tex,
% Reálná čísla a Číselná osa (bez konstrukcí dalších odmocnin a bez intervalů).
% Realisticky.cz, prostudované části v místní knihovně externi_zdroje/realisticky/:
% 01_Zakladni_poznatky/02_Ciselne_obory/010205_Realna_cisla_I/010205_Lekce.pdf
% (úvod, Př. 4–5: zaokrouhlování, přenos přes devítku);
% 01_Zakladni_poznatky/02_Ciselne_obory/010206_Realna_cisla_II/010206_Lekce.pdf
% (úvodní metodická poznámka, rozbor znamének).
% Konkrétní úlohy, formulace a řešení jsou vlastní tvorba. Platné cifry nezařazeny.
% Cílená revize: vlastní TikZ diagram, srovnávací tabulka a úlohy V1–V2
% podle zadání učitele. Další podklady v uvedené kapitole učebnice:
% Druhy číselných oborů a tabulky vlastností N, Z, Q;
% místní Realisticky 010205, Př. 1 a dodatek o neuzavřenosti iracionálních čísel,
% a 010201, úvodní přehled oborů a Zajímavé vlastnosti operací.
% Diagram: vnořené oblasti pro skutečné inkluze, nikoli velikosti nebo číselná osa.
\documentclass[11pt,a4paper]{article}
% Spolecny styl pro ucitelske pripravy MAT (Work/MAT)
% Verze 1.1, 2026-10-23
\NeedsTeXFormat{LaTeX2e}
\ProvidesPackage{priprava}[2026/10/23 v1.1 Sablona ucitelskych priprav MAT]

% Zakladni typografie a cestina
\RequirePackage[T1]{fontenc}
\RequirePackage[utf8]{inputenc}
\RequirePackage[czech]{babel}
\RequirePackage{lmodern}
\RequirePackage{microtype}
\RequirePackage{amsmath,amssymb}
\RequirePackage{xcolor}
\RequirePackage{enumitem}
\RequirePackage[most]{tcolorbox}
\RequirePackage{tabularx,booktabs,array}
\RequirePackage{fancyhdr}
\RequirePackage{needspace}
\RequirePackage{geometry}
\RequirePackage[hidelinks]{hyperref}

\geometry{a4paper,top=17mm,bottom=18mm,left=23mm,right=23mm}
\definecolor{notegray}{gray}{0.39}
\setlength{\parindent}{0pt}
\setlength{\parskip}{2.5pt plus 1pt minus 0.5pt}
\setlist[itemize]{topsep=2pt,itemsep=1pt,parsep=0pt,leftmargin=1.45em}
\setlist[enumerate]{topsep=2pt,itemsep=2pt,parsep=0pt,leftmargin=1.65em}
\setcounter{secnumdepth}{0}
\renewcommand{\arraystretch}{1.12}

% Zahlavi: vlevo rocnik | globalni tema, vpravo cislo: nazev hodiny
% Zapati zustava zcela prazdne.
\pagestyle{fancy}
\fancyhf{}
\renewcommand{\headrulewidth}{0.35pt}
\renewcommand{\footrulewidth}{0pt}
\setlength{\headheight}{14pt}
\newcommand{\nastavPripravu}[4]{%
  \fancyhead[L]{\small #1\enspace\textbar\enspace #2}%
  \fancyhead[R]{\small #3: #4}%
}

% Minimalisticke cernobile boxy. Kazda definice zvlast.
\tcbset{commonlesson/.style={enhanced,breakable=false,colback=white,colframe=black,
  colbacktitle=white,coltitle=black,boxrule=0.45pt,
  arc=0.8mm,outer arc=0.8mm,left=3mm,right=3mm,top=1.6mm,bottom=1.6mm,
  toptitle=1mm,bottomtitle=0.7mm,fonttitle=\bfseries,
  attach boxed title to top left={xshift=2.5mm,yshift=-1.15mm},
  boxed title style={colback=white,colframe=white,boxrule=0pt,left=1mm,right=1mm}}}
\newtcolorbox{zadaniBox}[1][]{commonlesson,title={Zadání},#1}
\newtcolorbox{prikladBox}[1][]{commonlesson,title={Příklad},#1}
\newtcolorbox{definiceBox}[1][]{commonlesson,title={Definice},#1}

% Spolecny box pro zadani a jeho reseni, bez mezery mezi castmi.
% Pouziti:
%   \begin{prikladReseniBox}
%     Zadani ulohy.
%     \tcblower
%     \begin{align*} ... \end{align*}
%   \end{prikladReseniBox}
% Spodni cast ma jemny sedy podklad a delici linku; lze v ni sazet
% vice radku matematiky. Obe casti maji jen jeden spolecny ramecek.
\tcbset{solutionlesson/.style={
  commonlesson,
  enhanced jigsaw,
  breakable,
  colbacklower=black!3,
  segmentation style={solid,draw=black!25,line width=0.35pt},
  fontlower=\small\color{notegray},
  before lower={\noindent\textbf{Řešení:}\enspace},
  leftlower=3mm,rightlower=3mm
}}
\newtcolorbox{prikladReseniBox}[1][]{solutionlesson,title={Příklad},#1}
\newtcolorbox{zadaniReseniBox}[1][]{solutionlesson,title={Zadání},#1}

% Pomucky pro autora pripravy: sedive poznamky; \reseni pro starsi pripravy
\newcommand{\poznamka}[1]{\par\begingroup\small\color{notegray}\itshape
  \textnormal{\textbf{Poznámka pro učitele.}}\enspace #1\par\endgroup}
\newcommand{\reseni}[1]{\par\begingroup\small\color{notegray}
  \textbf{Očekávané řešení:}\enspace #1\par\endgroup}
\newcommand{\blok}[1]{\needspace{5\baselineskip}\subsection*{#1}}
\newcommand{\vyrok}[1]{\uv{#1}}

\nastavPripravu{1. ročník}{Číselné obory}{07}{Reálná čísla a zaokrouhlování}
\begin{document}

\poznamka{Rozšířená příprava je zásobníkem; všechna procvičení se nemusejí
vejít do 45 minut. Orientační výběr: rozšíření oboru a osa 10 min,
diagram s příklady 8 min, uzavřenost a její meze 10 min,
zaokrouhlování 10 min, vybrané samostatné úlohy a shrnutí 7 min.
Původní D1 je asi na 15 min; rezervu vybírej podle potřeb třídy.
Nepředbíhej úpravy odmocnin ani vědecký zápis.}

\blok{Která čísla ještě na ose chybějí?}
Zlomek $\frac18$ má konečný zápis $0{,}125$, zlomek $\frac13$ nekonečný
periodický zápis $0{,}\overline3$. Nekonečný počet desetinných míst tedy
sám o sobě neznamená, že číslo není racionální.
Úhlopříčka jednotkového čtverce má podle známé Pythagorovy věty délku
$\sqrt2$. Tuto přesnou délku jako podíl celých čísel zapsat nelze.

\begin{definiceBox}
\textbf{Iracionální číslo} je reálné číslo, které nelze vyjádřit
zlomkem $\frac pq$ s celými $p,q$ a $q\ne0$.
Jeho desetinný rozvoj je nekonečný a není od žádného místa periodický.
\end{definiceBox}

\begin{definiceBox}
\textbf{Reálná čísla} zahrnují racionální i iracionální čísla.
Jejich obor značíme $\mathbb R$ a geometricky je představujeme všemi
body číselné osy bez mezer.
\end{definiceBox}

Příklady iracionálních čísel jsou $\sqrt2$ a $\pi$. Naproti tomu
$\sqrt9=3$ je přirozené číslo. Značka odmocniny sama o sobě nerozhoduje.
Všechna přirozená čísla jsou celá, všechna celá racionální a všechna
racionální reálná. Iracionální číslo je reálné, ale není racionální.

\poznamka{Iracionalitu $\sqrt2$ a $\pi$ zde sděl jako známou vlastnost,
ne jako závěr z displeje kalkulačky. Konečný výpis číslic nedokazuje
neperiodičnost nekonečného rozvoje. Úplný důkaz nyní není cílem hodiny.}

\poznamka{Volitelně lze iracionalitu $\sqrt2$ zdůvodnit sporem:
předpokládej $\sqrt2=p/q$ v základním tvaru s kladnými celými $p,q$.
Pak $p^2=2q^2$, tedy $p$ je sudé (druhá mocnina lichého čísla je lichá).
Po dosazení $p=2k$ vyjde $q^2=2k^2$, takže i $q$ je sudé.
To odporuje základnímu tvaru. Důkaz je pro zájemce; v základní části
stačí geometrická motivace jednotkovým čtvercem.}

\begin{definiceBox}
\textbf{Číselná osa} je přímka se zvoleným počátkem $0$, jednotkou délky
a kladným směrem. Každému reálnému číslu odpovídá právě jeden bod a naopak.
Při obvyklém směru doprava leží větší čísla více vpravo.
\end{definiceBox}

\begin{prikladReseniBox}[title={P1: přesná hodnota a poloha},breakable=false]
Seřaď vzestupně a přibližně umísti na osu
$-\frac32$, $-\sqrt2$, $0$, $\frac54$, $\sqrt2$.
Použij $1{,}41<\sqrt2<1{,}42$.
\tcblower
$-1{,}42<-\sqrt2<-1{,}41$, proto
\[
-\frac32<-\sqrt2<0<\frac54<\sqrt2.
\]
Hodnoty $-\frac32=-1{,}5$ a $\frac54=1{,}25$ jsou přesné.
Číslo $\sqrt2$ je na obrázku umístěno pouze přibližně.
\begin{center}
\begin{tikzpicture}[x=2.7cm,y=0.75cm]
\draw[->] (-2.1,0)--(2.12,0);
\foreach \x in {-2,-1,0,1,2}{\draw (\x,-0.08)--(\x,0.08) node[below=5pt] {$\x$};}
\fill (-1.5,0) circle (1.4pt) node[above=9pt] {$-\frac32$};
\fill (-1.4142,0) circle (1.4pt) node[below=22pt] {$-\sqrt2$};
\fill (1.25,0) circle (1.4pt) node[above=9pt] {$\frac54$};
\fill (1.4142,0) circle (1.4pt) node[below=22pt] {$\sqrt2$};
\end{tikzpicture}
\end{center}
\end{prikladReseniBox}

\newpage
\blok{Vennův diagram číselných oborů}
Zápis
\[
\mathbb N\subset\mathbb Z\subset\mathbb Q\subset\mathbb R
\]
čteme: každé číslo z menšího oboru patří také do následujících větších
oborů. Zde jsou všechny uvedené inkluze vlastní: větší obor obsahuje
i další čísla. Značení zatím slouží ke čtení obrázku;
obecné vztahy a operace mezi množinami budeme probírat později.

\begin{center}
\begin{tikzpicture}[x=1cm,y=1cm,font=\normalsize]
\draw[rounded corners=3mm,fill=black!4,line width=0.8pt] (0,0) rectangle (13.2,6.2);
\node[anchor=west] at (0.35,5.8) {$\mathbb R$};
\draw[fill=white,line width=0.8pt] (4.4,3.1) ellipse [x radius=4cm,y radius=2.5cm];
\node at (4.4,5.25) {$\mathbb Q$};
\draw[fill=black!8,line width=0.8pt] (3.8,2.8) ellipse [x radius=2.8cm,y radius=1.6cm];
\node at (3.6,4.0) {$\mathbb Z$};
\draw[fill=white,line width=0.8pt] (3.1,2.65) ellipse [x radius=1.5cm,y radius=0.85cm];
\node at (3.1,3.05) {$\mathbb N$};
\node at (3.1,2.45) {$1,\ 2,\ 5$};
\node at (5.3,2.8) {$0$};
\node at (3.8,1.5) {$-3$};
\node at (6.9,3.2) {$\frac12$};
\node at (4.5,4.75) {$-\frac74$};
\node[align=center] at (10.7,4.8) {$\mathbb R\setminus\mathbb Q$\\iracionální čísla};
\node at (9.7,3.3) {$\sqrt2$};
\node at (11.6,2.65) {$-\sqrt2$};
\node at (10.5,1.5) {$\pi$};
\end{tikzpicture}
\end{center}

Číslo napsané uvnitř $\mathbb N$ leží současně uvnitř $\mathbb Z$,
$\mathbb Q$ i $\mathbb R$. Oblast $\mathbb R\setminus\mathbb Q$ tvoří
všechna reálná čísla mimo $\mathbb Q$, tedy právě iracionální čísla.
Nula je mimo $\mathbb N$, ale uvnitř $\mathbb Z$.
Velikosti ploch ani vzdálenosti v diagramu nic neříkají o velikosti
čísel nebo o počtu prvků.

\begin{prikladReseniBox}[title={V1: přečti vnořené oblasti},breakable=false]
a) Do kterých ze čtyř oborů patří $5$ a $0$?
b) Patří $\frac12$ do $\mathbb Z$? Je $-\sqrt2$ racionální?
c) Kam bys připsal $\sqrt9$ a $0{,}5$?
\tcblower
a) $5$ patří do všech čtyř; $0$ do $\mathbb Z,\mathbb Q,\mathbb R$.
b) $\frac12\notin\mathbb Z$ a $-\sqrt2\notin\mathbb Q$.
c) $\sqrt9=3$ do nejvnitřnější oblasti $\mathbb N$;
$0{,}5=\frac12$ do $\mathbb Q$ mimo $\mathbb Z$.
Různé zápisy téhož čísla nepředstavují různé prvky.
\end{prikladReseniBox}

\poznamka{Nech žáky jednotlivé oblasti skutečně ukázat prstem
a u každého příkladu říci i všechny větší obory. Nečti obrázek jako
několik oddělených přihrádek. Vnoření vyjadřuje vztahy oborů,
nikoli relativní polohy jejich čísel na ose.}

\newpage
\blok{Rozšíření oboru a uzavřenost operací}
Z potřeby počítat předměty vycházíme z $\mathbb N$. Chceme-li vždy
odečítat, potřebujeme také nulu a záporná čísla: například $2-5=-3$.
Chceme-li dělit i při necelém podílu, přidáme zlomky: $1:2=\frac12$.
Rozšíření číselných oborů tak souvisí mimo jiné s operacemi,
které v předchozím oboru nebylo vždy možné provést bez jeho opuštění.

Tabulka říká, zda pro \emph{každou} přípustnou dvojici čísel z uvedeného
oboru zůstane výsledek ve stejném oboru. \uv{Ne} neznamená, že se to
nepodaří nikdy. U dělení zkoumáme pouze dvojice s nenulovým dělitelem.

\begin{center}
\begin{tabular}{lcccc}
\toprule
Obor čísel $a,b$ & $a+b$ & $a-b$ & $a\cdot b$ & $a:b$, $b\ne0$\\
\midrule
$\mathbb N$ & ano & ne & ano & ne\\
$\mathbb Z$ & ano & ano & ano & ne\\
$\mathbb Q$ & ano & ano & ano & ano\\
$\mathbb R$ & ano & ano & ano & ano\\
$\mathbb R\setminus\mathbb Q$ & ne & ne & ne & ne\\
\bottomrule
\end{tabular}
\end{center}

\textbf{Dělení nulou není definováno v žádném z těchto oborů.}
\uv{Ano} v posledním sloupci pro $\mathbb Q$ a $\mathbb R$ tuto výjimku
neruší. V posledním řádku porovnáváme samostatně iracionální čísla;
nejde o všechny reálné hodnoty.

\begin{prikladReseniBox}[title={V2: zdůvodni záporné odpovědi protipříkladem},breakable=false]
Ke každému \uv{ne} v tabulce najdi dvojici příslušného typu.
U dělení se drž nenulového dělitele.
\tcblower
V $\mathbb N$ například $3-3=0\notin\mathbb N$.
Pro dělení v $\mathbb N$ i $\mathbb Z$ stačí $1:2=\frac12$, které není celé.
Pro dvě iracionální čísla dostáváme například
\begin{align*}
\sqrt2+(-\sqrt2)&=0, & \sqrt2-\sqrt2&=0,\\
\sqrt2\cdot\sqrt2&=2, & \sqrt2:\sqrt2&=1.
\end{align*}
Všechny výsledky jsou racionální. Protipříklady vyvracejí uzavřenost,
nikoli možnost správně tyto výpočty provést v širším oboru $\mathbb R$.
\end{prikladReseniBox}

V $\mathbb Q$ jsou již součet, rozdíl, součin a podíl nenulovým číslem
racionální: při počítání se zlomky dostaneme opět podíl celých čísel
s nenulovým jmenovatelem. Přesto potřebujeme širší $\mathbb R$,
například pro délku $\sqrt2$ nebo pro řešení $x^2=2$.
Přechod z $\mathbb Q$ do $\mathbb R$ tedy není potřebný k uzavření
těchto čtyř operací; doplňuje další skutečné hodnoty, například některé odmocniny.

\poznamka{Jeden příklad dokazuje odpověď \uv{ne}, ale ne obecné \uv{ano}.
U racionálních čísel se opři o obecná pravidla počítání se zlomky,
nikoli jen o několik dosazení. Neměň tabulku uzavřenosti na tabulku
komutativity či asociativity: odpovídá na jinou otázku.}

\newpage
\blok{Zaokrouhlení mění hodnotu zápisu}
\begin{definiceBox}
\textbf{Zaokrouhlení na daný řád} znamená nahrazení čísla nejbližším
násobkem jednotky tohoto řádu. Při přesné polovině zde volíme násobek
dále od nuly, například $2{,}45\doteq2{,}5$ a $-2{,}45\doteq-2{,}5$ na desetiny.
\end{definiceBox}

U kladného čísla sledujeme první vynechanou číslici: je-li menší než $5$,
ponechanou část neměníme; je-li alespoň $5$, přičteme jednotku posledního
ponechaného řádu. U záporného čísla stejně zaokrouhlíme jeho kladný
protějšek a vrátíme znaménko. Používáme $\doteq$, pokud zaokrouhlení
není přesná rovnost; příslušný řád vždy uvedeme.

Zaokrouhlením na setiny se od přesné hodnoty odchýlíme nejvýše
o $0{,}005$, tedy o polovinu kroku $0{,}01$ mezi sousedními setinami.
Například zápis $\sqrt2\doteq1{,}41$ neoznačuje nové přesné vyjádření
odmocniny. V mezivýpočtech proto ponecháváme přesný zlomek či odmocninu
a zaokrouhlujeme až požadovaný výsledek.

\begin{prikladReseniBox}[title={P2: devítky a záporné číslo},breakable=false]
Zaokrouhli na setiny: $3{,}146$, $9{,}995$, $-2{,}345$ a $-2{,}341$.
\tcblower
$3{,}146\doteq3{,}15$; $9{,}995\doteq10{,}00$;
$-2{,}345\doteq-2{,}35$; $-2{,}341\doteq-2{,}34$.
V druhém případě se přenos dotkne i celé části. Koncové nuly
v $10{,}00$ sdělují požadované zaokrouhlení na setiny.
\end{prikladReseniBox}

\begin{prikladReseniBox}[title={P3: rozhoduj z původního čísla},breakable=false]
Zaokrouhli $1{,}249$ přímo na desetiny. Porovnej s postupem,
při kterém se nejprve zaokrouhlí na setiny a pak znovu na desetiny.
\tcblower
Přímo: první vynechaná číslice je $4$, tedy $1{,}249\doteq1{,}2$.
Přes mezivýsledek: $1{,}249\doteq1{,}25\doteq1{,}3$.
Dvojí zaokrouhlování dalo jiný výsledek. Pro požadovaný řád proto
vycházíme z původního čísla.
\end{prikladReseniBox}

\blok{Samostatné procvičování}
\begin{zadaniReseniBox}[title={S1: obor, osa, přesnost},breakable=false]
\begin{enumerate}[label=\alph*)]
\item U každého čísla uveď všechny vhodné obory z $\mathbb N,\mathbb Z,
\mathbb Q,\mathbb R$; není-li racionální, řekni to:
$0$, $-\frac72$, $\sqrt{16}$, $0{,}\overline{27}$, $-\sqrt2$.
\item Seřaď a zakresli přibližně na osu: $-1{,}5$, $-\frac54$, $-\sqrt2$, $-1$.
\item Na setiny zaokrouhli $4{,}736$, $0{,}0049$, $-7{,}895$, $12{,}999$.
\item Na stovky zaokrouhli $14\,950$ a $-2\,449$.
\item Platí $\sqrt2=1{,}41$? Stačí zápis $0{,}123456\ldots$ k rozhodnutí,
zda jde o racionální číslo?
\end{enumerate}
\tcblower
a) $0$: $\mathbb Z,\mathbb Q,\mathbb R$;
$-\frac72$: $\mathbb Q,\mathbb R$;
$\sqrt{16}=4$: všechny čtyři;
$0{,}\overline{27}=\frac{27}{99}=\frac3{11}$: $\mathbb Q,\mathbb R$;
$-\sqrt2$: jen $\mathbb R$, je iracionální.
b) $-1{,}5<-\sqrt2<-\frac54<-1$.
c) $4{,}74$; $0{,}00$; $-7{,}90$; $13{,}00$.
d) $15\,000$; $-2\,400$.
e) Ne, pouze $\sqrt2\doteq1{,}41$ na setiny.
Konečný začátek rozvoje nerozhoduje o jeho dalším průběhu.
\end{zadaniReseniBox}

\poznamka{Při zakreslování trvej na stejné délce jednotky v kladné
i záporné části osy. Závěr: každý žák rozliší přesné $\frac14=0{,}25$
od přibližného $\pi\doteq3{,}14$. Příště budeme vyznačovat celé úseky osy.}

\newpage
\blok{Rezerva: proč výsledek platí}
\begin{zadaniReseniBox}[title={R1: volitelně, 10–15 minut},breakable=false]
\begin{enumerate}[label=\alph*)]
\item Ověř $1{,}41<\sqrt2<1{,}42$ porovnáním druhých mocnin.
Vysvětli, proč je možné takto porovnávat právě kladná čísla.
\item Může být součet dvou iracionálních čísel racionální?
Uveď příklad. Dokaž, že $1+\sqrt2$ je iracionální.
\item Změřená kladná délka se po zaokrouhlení na desetiny zobrazila
jako $2{,}4$ m. Popiš slovy a nerovností všechny možné původní délky.
\end{enumerate}
\tcblower
a) $1{,}41^2=1{,}9881<2<2{,}0164=1{,}42^2$.
U nezáporných čísel zachovává druhá mocnina uspořádání; u záporných
čísel to obecně neplatí, například $-3<-2$, ale $9>4$.
b) Ano: $\sqrt2+(-\sqrt2)=0$. Kdyby $1+\sqrt2$ bylo racionální,
odečtením racionální jedničky by bylo racionální i $\sqrt2$, což není.
c) $2{,}35\le d<2{,}45$ v metrech. Dolní mez se zaokrouhlí na $2{,}4$,
horní už na $2{,}5$. Nejde o jedinou přesnou původní délku.
\end{zadaniReseniBox}

\blok{Domácí procvičování}
\begin{zadaniReseniBox}[title={D1: přibližně 15 minut},breakable=false]
\begin{enumerate}[label=\arabic*)]
\item Uveď všechny vhodné obory $\mathbb N,\mathbb Z,\mathbb Q,\mathbb R$
pro $\sqrt{25}$, $-6$, $\frac29$ a $\pi$.
\item Seřaď a přibližně vyznač na ose $-\frac74$, $-\sqrt2$, $0$, $\frac32$.
\item Na desetiny zaokrouhli $6{,}24$, $6{,}25$, $-6{,}25$, $9{,}96$.
\item Na tisíciny zaokrouhli $0{,}07856$ a $-3{,}0005$.
\item Číslo $2{,}449$ zaokrouhli přímo na desetiny a pak porovnej s postupem
přes setiny. Který výsledek odpovídá zadání?
\item Napiš racionální a iracionální číslo ostře mezi $1$ a $2$.
Vysvětli, proč konečný údaj $1{,}4142$ není totéž co $\sqrt2$.
\end{enumerate}
\tcblower
1. $\sqrt{25}=5$: všechny čtyři; $-6$: $\mathbb Z,\mathbb Q,\mathbb R$;
$\frac29$: $\mathbb Q,\mathbb R$; $\pi$: pouze $\mathbb R$, iracionální.
2. $-\frac74<-\sqrt2<0<\frac32$.
3. $6{,}2$; $6{,}3$; $-6{,}3$; $10{,}0$.
4. $0{,}079$; $-3{,}001$.
5. Přímo $2{,}4$, přes $2{,}45$ dostaneme $2{,}5$; správný je přímý výsledek.
6. Například $\frac32$ a $\sqrt2$. Konečné $1{,}4142=\frac{7071}{5000}$
je racionální, $\sqrt2$ iracionální; jde pouze o přiblížení.
\end{zadaniReseniBox}

\end{document}
